\documentclass[12pt, a4paper]{ctexart}
\usepackage{amsmath, amssymb, amsthm}
\usepackage{fancyhdr}
\usepackage{geometry}
\usepackage{enumitem}
\usepackage{graphicx}
\usepackage{tikz}
\usepackage{pgfplots}
\pgfplotsset{compat=1.18}

\geometry{left=2cm, right=2cm, top=2.5cm, bottom=2.5cm}

\setCJKmainfont{Hiragino Sans GB}
\setCJKsansfont{Heiti SC}

\pagestyle{fancy}
\fancyhf{}
\fancyhead[L]{高考数学满分突破·复数与欧拉公式}
\fancyhead[R]{测试卷}
\fancyfoot[C]{\thepage}
\renewcommand{\headrulewidth}{0.4pt}

\setlength{\parindent}{0pt}
\setlength{\parskip}{6pt}

\newcommand{\score}[1]{\hfill\textbf{（#1分）}}

\begin{document}

\begin{center}
{\LARGE\bfseries 高考数学满分突破测试卷（九）}\\[8pt]
{\large 复数与欧拉公式}\\[6pt]
考试时间：120分钟 \quad 满分：150分
\end{center}

\vspace{10pt}

{\large\bfseries 一、单项选择题：本题共8小题，每小题5分，共40分。每小题给出的四个选项中，只有一项是符合题目要求的。}

\vspace{8pt}

\begin{enumerate}

\item 设 $z=\dfrac{1+i}{1-i}$，则 $|z|=$ \score{5}
\begin{enumerate}[label=\Alph*.]
    \item $0$
    \item $1$
    \item $\sqrt{2}$
    \item $2$
\end{enumerate}

\item 设 $z_1=1+2i$，$z_2=3-i$，则 $z_1\cdot z_2=$ \score{5}
\begin{enumerate}[label=\Alph*.]
    \item $5+5i$
    \item $5-5i$
    \item $1+5i$
    \item $3+5i$
\end{enumerate}

\item 由欧拉公式 $e^{i\theta}=\cos\theta+i\sin\theta$，则 $e^{i\pi}+1=$ \score{5}
\begin{enumerate}[label=\Alph*.]
    \item $-2$
    \item $-1$
    \item $0$
    \item $1$
\end{enumerate}

\item 设 $z=3+4i$，则 $\dfrac{1}{z}=$ \score{5}
\begin{enumerate}[label=\Alph*.]
    \item $\dfrac{3}{25}-\dfrac{4}{25}i$
    \item $\dfrac{3}{25}+\dfrac{4}{25}i$
    \item $\dfrac{4}{25}-\dfrac{3}{25}i$
    \item $\dfrac{3}{7}-\dfrac{4}{7}i$
\end{enumerate}

\item $1$ 的 $n$ 次单位根为 $\omega_k=e^{i\cdot\frac{2k\pi}{n}}$，$k=0,1,\cdots,n-1$。当 $n=3$ 时，$\omega_1^2+\omega_1+1=$ \score{5}
\begin{enumerate}[label=\Alph*.]
    \item $-1$
    \item $0$
    \item $1$
    \item $3$
\end{enumerate}

\item 设 $z=\cos\dfrac{\pi}{6}+i\sin\dfrac{\pi}{6}$，则 $z^6=$ \score{5}
\begin{enumerate}[label=\Alph*.]
    \item $1$
    \item $-1$
    \item $i$
    \item $-i$
\end{enumerate}

\item 用De Moivre公式，$(\cos\theta+i\sin\theta)^3=\cos 3\theta+i\sin 3\theta$，展开后取实部可得 $\cos 3\theta=$ \score{5}
\begin{enumerate}[label=\Alph*.]
    \item $4\cos^3\theta-3\cos\theta$
    \item $3\cos\theta-4\cos^3\theta$
    \item $4\cos^3\theta+3\cos\theta$
    \item $2\cos^3\theta-\cos\theta$
\end{enumerate}

\item 设 $z=re^{i\theta}$（$r>0$，$\theta\in\mathbb{R}$），则 $\bar{z}=$ \score{5}
\begin{enumerate}[label=\Alph*.]
    \item $re^{-i\theta}$
    \item $-re^{i\theta}$
    \item $re^{i(\pi-\theta)}$
    \item $re^{i(-\pi+\theta)}$
\end{enumerate}

\end{enumerate}

\vspace{10pt}

{\large\bfseries 二、多项选择题：本题共4小题，每小题5分，共20分。每小题有多个选项符合题目要求，全部选对得5分，部分选对得2分，有选错的得0分。}

\vspace{8pt}

\begin{enumerate}[resume]

\item 关于复数的基本性质，下列正确的是 \score{5}
\begin{enumerate}[label=\Alph*.]
    \item $|z_1z_2|=|z_1|\cdot|z_2|$
    \item $\overline{z_1+z_2}=\bar{z}_1+\bar{z}_2$
    \item $\overline{z_1z_2}=\bar{z}_1\cdot\bar{z}_2$
    \item $|z|^2=z\cdot\bar{z}$
\end{enumerate}

\item 关于欧拉公式及其推论，下列正确的是 \score{5}
\begin{enumerate}[label=\Alph*.]
    \item $e^{i\theta}=\cos\theta+i\sin\theta$
    \item $\cos\theta=\dfrac{e^{i\theta}+e^{-i\theta}}{2}$
    \item $\sin\theta=\dfrac{e^{i\theta}-e^{-i\theta}}{2i}$
    \item $|e^{i\theta}|=1$ 对所有 $\theta\in\mathbb{R}$ 成立
\end{enumerate}

\item 关于单位根，下列正确的是 \score{5}
\begin{enumerate}[label=\Alph*.]
    \item $x^n-1=\displaystyle\prod_{k=0}^{n-1}(x-\omega_k)$，其中 $\omega_k=e^{i\frac{2k\pi}{n}}$
    \item $\displaystyle\sum_{k=0}^{n-1}\omega_k^m=\begin{cases}n,&n\mid m\\0,&n\nmid m\end{cases}$
    \item 所有 $n$ 次单位根在复平面上构成正 $n$ 边形的顶点
    \item $\omega_k^n=1$ 对所有 $k$ 成立
\end{enumerate}

\item 用复数方法求三角级数的和，下列正确的是 \score{5}
\begin{enumerate}[label=\Alph*.]
    \item $\displaystyle\sum_{k=0}^{n-1}\cos k\theta=\dfrac{\sin\frac{n\theta}{2}}{\sin\frac{\theta}{2}}\cos\dfrac{(n-1)\theta}{2}$（$\theta\neq 2m\pi$）
    \item $\displaystyle\sum_{k=0}^{n-1}\sin k\theta=\dfrac{\sin\frac{n\theta}{2}}{\sin\frac{\theta}{2}}\sin\dfrac{(n-1)\theta}{2}$
    \item 上述求和可以通过对等比数列 $\displaystyle\sum_{k=0}^{n-1}e^{ik\theta}$ 取实部和虚部得到
    \item 当 $\theta=\dfrac{2\pi}{n}$ 时，$\displaystyle\sum_{k=0}^{n-1}\cos k\theta=0$
\end{enumerate}

\end{enumerate}

\vspace{10pt}

{\large\bfseries 三、填空题：本题共4小题，每小题5分，共20分。}

\vspace{8pt}

\begin{enumerate}[resume]

\item 设 $z=\sqrt{3}+i$，则 $z$ 的模 $|z|=$ \underline{\hspace{2cm}}，辐角主值 $\arg z=$ \underline{\hspace{2cm}}。\score{5}

\item 设 $z_1=2e^{i\frac{\pi}{3}}$，$z_2=3e^{i\frac{\pi}{6}}$，则 $z_1z_2=$ \underline{\hspace{2cm}}，$\dfrac{z_1}{z_2}=$ \underline{\hspace{2cm}}（用指数形式表示）。\score{5}

\item 利用欧拉公式求和 $\displaystyle\sum_{k=0}^{4}\cos\dfrac{2k\pi}{5}=$ \underline{\hspace{2cm}}。\score{5}

\item 方程 $z^4=-1$ 在复数范围内的所有解为 \underline{\hspace{2cm}}。\score{5}

\end{enumerate}

\vspace{10pt}

{\large\bfseries 四、解答题：本题共6小题，共70分。解答应写出文字说明、证明过程或演算步骤。}

\vspace{8pt}

\begin{enumerate}[resume]

\item （10分）从复数的几何意义出发。

\begin{enumerate}[label=(\arabic*)]
    \item 证明复数乘法的几何意义：设 $z_1=r_1e^{i\theta_1}$，$z_2=r_2e^{i\theta_2}$，则 $z_1z_2$ 对应将 $z_1$ 的模放大 $r_2$ 倍，辐角增加 $\theta_2$；
    \item 在复平面上，设 $A$ 对应 $z_1=1+i$，$B$ 对应 $z_2=2+i$，求三角形 $OAB$（$O$ 为原点）的面积；
    \item 证明：$|z_1+z_2|\leq |z_1|+|z_2|$（复数的三角不等式）。
\end{enumerate}

\item （12分）用欧拉公式推导三角恒等式。

\begin{enumerate}[label=(\arabic*)]
    \item 利用 $e^{i(\alpha+\beta)}=e^{i\alpha}\cdot e^{i\beta}$ 证明和角公式 $\cos(\alpha+\beta)=\cos\alpha\cos\beta-\sin\alpha\sin\beta$；
    \item 利用De Moivre公式证明 $\cos 4\theta=8\cos^4\theta-8\cos^2\theta+1$；
    \item 求 $\cos\dfrac{\pi}{5}$ 的精确值（提示：利用 $z^5=1$ 的根）。
\end{enumerate}

\item （12分）复数与多项式。

\begin{enumerate}[label=(\arabic*)]
    \item 证明代数基本定理的推论：$n$ 次复系数多项式在复数范围内恰有 $n$ 个根（计重数）；
    \item 设 $p(z)=z^4+z^3+z^2+z+1$，利用 $p(z)=\dfrac{z^5-1}{z-1}$（$z\neq 1$）求 $p(z)$ 的所有根；
    \item 利用上述结果证明 $\cos\dfrac{2\pi}{5}+\cos\dfrac{4\pi}{5}=-\dfrac{1}{2}$。
\end{enumerate}

\item （12分）用复数方法求解三角求和问题。

\begin{enumerate}[label=(\arabic*)]
    \item 设 $S=\displaystyle\sum_{k=0}^{n-1}\cos k\theta$，$T=\displaystyle\sum_{k=0}^{n-1}\sin k\theta$。利用 $S+iT=\displaystyle\sum_{k=0}^{n-1}e^{ik\theta}$ 求 $S$ 和 $T$（$\theta\neq 2m\pi$）；
    \item 利用上述结果求 $\displaystyle\sum_{k=0}^{9}\cos\dfrac{k\pi}{5}$ 和 $\displaystyle\sum_{k=0}^{9}\sin\dfrac{k\pi}{5}$；
    \item 求 $\displaystyle\sum_{k=1}^{n}k\cos k\theta$（$\theta\neq 2m\pi$）。
\end{enumerate}

\item （12分）复数在高考解题中的高观点应用。

\begin{enumerate}[label=(\arabic*)]
    \item 设 $\omega$ 为 $1$ 的三次单位根（$\omega\neq 1$），证明 $1+\omega+\omega^2=0$；
    \item 利用上述结论求 $\left(1-\omega+\omega^2\right)^6$ 的值；
    \item 求 $\displaystyle\sum_{k=0}^{n}\dbinom{n}{k}\cos\dfrac{2k\pi}{3}$ 的值（提示：考虑 $(1+e^{i\frac{2\pi}{3}})^n$ 的实部）。
\end{enumerate}

\item （12分）复数与线性变换。

\begin{enumerate}[label=(\arabic*)]
    \item 证明复数乘法 $w=az$（$a\neq 0$ 为复常数）在复平面上对应一个旋转变换与伸缩变换的复合；
    \item 设变换 $w=\dfrac{z-i}{z+i}$，求 $w$ 将实轴映射到什么曲线（提示：设 $z=x\in\mathbb{R}$，求 $|w|$）；
    \item 证明：上述变换将上半平面 $\text{Im}(z)>0$ 映射到单位圆 $|w|<1$ 的内部。
\end{enumerate}

\end{enumerate}

\end{document}
